\chapter{Limitations and Honest Boundaries}
\label{chap:limitations}

\section{Native feature limits}

\begin{table}[htbp]
\centering
\caption{Current limitations of the native implementation and why they matter.}
\label{tab:limitations}
\begin{tabularx}{\textwidth}{@{}p{0.26\textwidth} p{0.22\textwidth} X@{}}
\toprule
Limitation & Status & Consequence \\
\midrule
Head dimension cap & $D \le 256$ & keeps register and shared-memory usage simple, but excludes larger heads \\
No dropout & unsupported & not a drop-in training replacement for general Transformer workloads \\
No arbitrary mask or bias & unsupported & no padding masks, additive biases, ALiBi, or custom sparsity patterns \\
No GQA/MQA & unsupported & Q, K, and V must have the same head count \\
No second-order gradients & unsupported & backward is sufficient for ordinary training but not higher-order differentiation \\
No tensor-core path & omitted by design & likely leaves substantial performance on the table \\
Fixed tile sizes & fixed at compile time & easy to analyze, but not tuned per architecture or workload \\
Authoring environment has no GPU & currently true & empirical claims must wait for Kaggle or another CUDA machine \\
\bottomrule
\end{tabularx}
\end{table}

\section{Validation limits}

The present tests are strong enough to support first-order correctness claims on the exercised grid, but they are not exhaustive. In particular:
\begin{itemize}
  \item a passing suite does not rule out all race conditions or OOB accesses;
  \item the suite does not yet cover every possible GPU architecture or driver stack;
  \item numerical tolerance checks are necessary but cannot prove ideal conditioning;
  \item benchmark-based memory measurements rely on allocator statistics rather than a complete device-memory trace.
\end{itemize}

\section{Performance interpretation limits}

Any future performance discussion in this repository must respect three caveats.

\paragraph{First,} a dense-equivalent TFLOP/s estimate is not a true instruction-throughput metric. It is a convenient normalization based on the mathematical work model.

\paragraph{Second,} benchmark results on Kaggle's GPU should not be generalized automatically to other architectures. Shared-memory capacity, tensor-core behavior, cache geometry, and compiler quality all matter.

\paragraph{Third,} comparison with PyTorch SDPA is comparison with a moving target. PyTorch may route to math, efficient, cuDNN, or FlashAttention-like backends depending on version, dtype, mask, and hardware \citep{pytorch212sdpa}. That is why the benchmark records backend-related metadata and allows explicit SDPA-backend forcing.

\section{Scientific scope}

This project should be read as a serious implementation study, not as a claim to state-of-the-art attention performance. Its main scientific scope is narrower and more defensible:
\begin{enumerate}
  \item demonstrate the exact online-softmax algorithm in code;
  \item implement the corresponding first-order backward pass;
  \item expose the operator through a modern PyTorch extension interface;
  \item define a reproducible benchmark and reporting workflow.
\end{enumerate}

\section{Future work}

Natural next steps are clear:
\begin{itemize}
  \item tune tile sizes by head dimension and architecture;
  \item add vectorized loads and tensor-core kernels;
  \item support dropout, masks, and head-count asymmetry;
  \item add profiler-driven optimization and sanitization passes;
  \item automate more report generation directly from benchmark and profiler artifacts.
\end{itemize}

The value of the current repository is that those directions are now grounded in a complete, understandable baseline rather than in an unexplained black box.
